% ---------------------------------------------------------------------
\subsection{Processamento descentralizado}\label{subsec:dbp}
% ---------------------------------------------------------------------

Na arquitetura descentralizada, as $M$ antenas da \sout{estação base} \textcolor{red}{BS} são particionadas em $C$ clusters, cada um associado a uma DU \textcolor{red}{(\textit{Distributed Unit}), uma unidade de processamento local \cite{zhao2025jsac}}. O vetor de ruído pode, então, ser organizado em componentes locais:

\[
\mathbf n=
\begin{bmatrix}
\mathbf n_1\\
\mathbf n_2\\
\vdots\\
\mathbf n_C
\end{bmatrix}.
\]

Essa partição induz uma estrutura em blocos na matriz de covariância $\mathbf R_n$. Os blocos diagonais $\mathbf R_{cc}$ descrevem a covariância do ruído observado pelas antenas do cluster $c$, enquanto os blocos fora da diagonal, $\mathbf R_{mn}$, representam as correlações entre antenas pertencentes a clusters distintos.

Cada DU dispõe localmente das amostras de ruído de suas próprias antenas. Assim, a DU $c$ pode estimar seu bloco diagonal sem trocar dados com as demais unidades:

\[
\hat{\mathbf R}_{cc}=
\frac{1}{N}
\sum_{i=1}^{N}
\mathbf n_c^i(\mathbf n_c^i)^H.
\]

Os blocos fora da diagonal, por outro lado, dependem de amostras provenientes de duas ou mais DUs. Sua estimação direta exige a troca de informações entre as unidades, aumentando o tráfego de comunicação.

Essa limitação estabelece uma relação entre as propriedades estatísticas do ruído e a estratégia de equalização descentralizada. 
